% TEMPLATE-MILC-v3.tex - plantilla maestra UNICA de las guias MILC v3.
%
% La usan las cuatro familias de guias, cada una con su driver:
%   · Grados 6-11          -> scripts/build-guias-g11.py
%   · Bebras               -> scripts/build-guias-bebras.py
%   · Semillero            -> scripts/build-guias-semillero.py
%   · Territorio interior  -> scripts/build-guias-territorio-interior.py
%
% Antes habia tres copias casi identicas (~400 lineas iguales, 6 distintas):
% cada mejora habia que aplicarla tres veces, y en la practica no se hacia
% (el motor de recursos vivia solo en dos de ellas). Lo que varia entre
% familias esta parametrizado en 6 placeholders que cada driver llena
% (se nombran aqui SIN su sintaxis real: el builder reemplaza por texto plano,
% tambien dentro de los comentarios, y un valor multilinea rompe el preambulo):
%
%   HEADER_IZQ        encabezado izquierdo de pagina
%   HEADER_DER        encabezado derecho de pagina
%   PORTADA_ETIQUETA  rotulo sobre el numero grande (GRADO/MOMENTO/MODULO)
%   PORTADA_SERIE     linea de serie de la portada
%   PORTADA_ID        identificador de la guia en la portada
%   PORTADA_CIRCULO   el circulito de grado (°), o vacio si no aplica
%   PDF_TITULO        titulo en texto plano (sin \textbf ni comillas TeX) para
%                     los metadatos del PDF (/Title); lo produce lib_xelatex.texto_pdf
%
% Composicion (auditoria 2026-09-03): las bandas de estacion piden espacio
% (\Needspace*) y prohiben el salto justo despues (\nopagebreak), asi no quedan
% huerfanas al pie; las cajas largas (softbox, drawbox, infoband, puentebox) son
% `breakable`, asi una caja de media pagina ya no arrastra todo a la siguiente
% dejando la anterior al 40 %. Todo texto sobre mostaza o verde va en milcNegro
% (blanco sobre #E5B400 da 1,93:1 y sobre #58A923 2,95:1; ninguno pasa WCAG AA).
% El resto de claves las documenta content/guias/_SCHEMA.md. El script valida
% que no queden placeholders sin reemplazar antes de escribir el archivo final.

% Etiquetado PDF (latex-lab): /Lang, estructura y texto alternativo de las
% figuras, para lectores de pantalla. Debe ir ANTES de \documentclass.
\DocumentMetadata{lang=es-CO,tagging=on}
\documentclass[11pt,letterpaper]{article}
% headheight=24pt: la cabecera derecha lleva dos lineas (identidad de la guia
% y estacion actual). headsep se acorta en la misma medida para que el bloque
% de texto quede exactamente donde estaba con headheight=12pt.
\usepackage[margin=2.0cm,headheight=24pt,headsep=13pt]{geometry}
\usepackage[table]{xcolor}
\usepackage{fontspec}
\usepackage[spanish]{babel}
\usepackage{tikz}
% Librerías TikZ para los diagramas de las guías (bloque `recursos:`):
%  · babel  → babel en español vuelve activos caracteres como «>», y TikZ los usa
%             en sus opciones (>=stealth, ->). Sin esta librería, cualquier
%             diagrama con flechas revienta con «\language@active@arg> has an
%             extra }». Debe ir ANTES de que se dibuje nada.
%  · positioning        → sintaxis «below left=2pt and 3pt».
%  · decorations.pathreplacing → llaves ({brace}) para señalar rangos.
\usetikzlibrary{babel,positioning,decorations.pathreplacing}
\usepackage{graphicx}
% Circuitos: el trabajo de electrónica se hace hoy con micro:bit y sus sensores
% integrados (MakeCode, sin cableado), así que no se necesitan esquemáticos.
% Si algún día entran montajes con protoboard, basta descomentar esta línea:
% los diagramas .tex/.tikz declarados en `recursos:` ya se incrustan solos.
% \usepackage{circuitikz}
\usepackage[most]{tcolorbox}
\usepackage{tabularx}
\usepackage{array}
\usepackage{fancyhdr}
\usepackage{titlesec}
\usepackage[document]{ragged2e}  % bandera a la derecha en todo el cuerpo
\usepackage{enumitem}
\usepackage{microtype}
\usepackage{etoolbox}
\usepackage{needspace}
% hyperref al final del preambulo: metadatos del PDF (titulo, autor, idioma,
% familia), marcadores por estacion (\pdfbookmark) y enlaces sin recuadro.
\usepackage[hidelinks,
  pdftitle={Sustentación pública y portafolio digital del grado 8},
  pdfauthor={PhD. Álvaro Cárdenas Orozco},
  pdfsubject={Tecnología e Informática · Grado 8 · Guía 30 · MILC v3},
  pdflang={es-CO},
  bookmarksopen=true,bookmarksnumbered=false]{hyperref}

% Helvetica Neue no trae la flecha U+2192, asi que cada «→» escrito literalmente
% en el YAML se compilaba a NADA, en silencio (462 flechas en 61 guias: rutas del
% tipo «paleta → tipografia → lockup» salian rotas). Mapeamos el caracter a su
% version matematica, que si tiene glifo. Asi el autor puede seguir escribiendo
% «→» comodamente en el YAML.
\usepackage{newunicodechar}
\newunicodechar{→}{\ensuremath{\rightarrow}}
% Mismo problema con los simbolos de marca y vinetas: el «✓» de «una palomita ✓»
% salia como cuadrito vacio en el PDF de todas las guias que lo usaban.
\usepackage{pifont}
\newunicodechar{✓}{\ding{51}}
\newunicodechar{✗}{\ding{55}}
\newunicodechar{✔}{\ding{52}}
\newunicodechar{•}{\textbullet}
\newunicodechar{≥}{\ensuremath{\geq}}
\newunicodechar{≤}{\ensuremath{\leq}}

\setmainfont{Helvetica Neue}
\setsansfont{Helvetica Neue}
\setlength{\parindent}{0pt}
\setlength{\parskip}{6pt}
% Sin particion de palabras: en 6.º-7.º una palabra cortada por guion al final
% de linea («Comuni-cacion») frena la lectura mucho mas que una linea corta.
\hyphenpenalty=10000
\exhyphenpenalty=10000
\pretolerance=2000
\tolerance=3000
\setlength{\emergencystretch}{3em}
\linespread{1.10}
\renewcommand{\arraystretch}{1.25}
\setlist{leftmargin=5mm,itemsep=1mm,topsep=1mm}
\newcolumntype{Y}{>{\RaggedRight\arraybackslash}X}

\definecolor{milcMagenta}{HTML}{E6007E}
\definecolor{milcVino}{HTML}{5A0038}
\definecolor{milcTurquesa}{HTML}{0093A5}
\definecolor{milcVerde}{HTML}{58A923}
\definecolor{milcMostaza}{HTML}{E5B400}
\definecolor{milcOcre}{HTML}{B86600}
\definecolor{milcNegro}{HTML}{241816}
\definecolor{milcGris}{HTML}{F7F7F4}
\definecolor{milcLinea}{HTML}{E8E2DD}
\definecolor{milcGradePrimary}{HTML}{FF6600}
\definecolor{milcGradeDark}{HTML}{9A3E00}
\definecolor{milcGradeSoft}{HTML}{FFE4D0}
\definecolor{milcGradeAccent}{HTML}{CFE7FF}
\definecolor{milcGradeText}{HTML}{FFFFFF}
\color{milcNegro}

\pagestyle{fancy}
\fancyhf{}
% Cabecera de dos lineas: identidad de la guia arriba y, a la derecha y en
% gris, la estacion en curso (\markright desde \milcestacion). Antes de la
% primera estacion la segunda linea va vacia; el \strut mantiene la altura
% para que la linea de arriba no baile entre paginas.
\lhead{\small Tecnología e Informática · Grado 8\\[-1pt]{\footnotesize\strut}}
\rhead{\small Guía 30 · MILC v3\\[-1pt]{\footnotesize\strut\color{milcNegro!65}\rightmark}}
\cfoot{\small\thepage}
\renewcommand{\headrulewidth}{0pt}

% Sobre mostaza (#E5B400) y verde (#58A923) el texto va en milcNegro; sobre el
% resto de la paleta, en blanco. Se usa dentro de fonttitle/fontupper, donde un
% \color posterior manda sobre coltitle.
\newcommand{\milctintasobre}[1]{%
  \ifboolexpr{test{\ifstrequal{#1}{milcMostaza}} or test{\ifstrequal{#1}{milcVerde}}}%
    {\color{milcNegro}}{\color{white}}}

\newtcolorbox{titlebox}[1]{
  enhanced,colback=#1,colframe=#1,arc=7mm,boxrule=0pt,
  left=7mm,right=7mm,top=3mm,bottom=3mm,
  before skip=8pt,after skip=8pt,
  fontupper=\sffamily\bfseries\Large\color{white}
}

\newtcolorbox{softbox}[3]{
  enhanced,breakable,pad at break=3mm,
  colback=#2,colframe=#1,borderline west={4pt}{0pt}{#1},
  arc=2mm,boxrule=.4pt,left=4mm,right=4mm,top=3mm,bottom=3mm,
  before skip=6pt,after skip=8pt,title={#3},coltitle=white,
  colbacktitle=#1,fonttitle=\sffamily\bfseries\milctintasobre{#1},fontupper=\RaggedRight,
  attach boxed title to top left={xshift=3mm,yshift=-2mm},
  boxed title style={arc=2mm, boxrule=0pt}
}

\newtcolorbox{drawbox}[2]{
  enhanced,breakable,pad at break=4mm,
  colback=white,colframe=#1,arc=4mm,boxrule=1pt,
  left=5mm,right=5mm,top=4mm,bottom=4mm,before skip=8pt,after skip=8pt,
  title={#2},coltitle=white,colbacktitle=#1,fonttitle=\sffamily\bfseries\milctintasobre{#1},
  fontupper=\RaggedRight,
  attach boxed title to top left={xshift=4mm,yshift=-2mm},
  boxed title style={arc=3mm, boxrule=0pt}
}

\newtcolorbox{infoband}[1]{
  enhanced,breakable,pad at break=2mm,colback=#1!8,colframe=#1!50,arc=2mm,boxrule=.5pt,
  left=4mm,right=4mm,top=2mm,bottom=2mm,before skip=4pt,after skip=4pt,
  fontupper=\small\RaggedRight
}

% Puente narrativo entre secciones (contrato editorial Conéctate).
% Da continuidad: "Acabas de X. Ahora vas a Y porque Z."
% Visualmente: banda izquierda en color del grado, texto en itálica pequeña.
\newtcolorbox{puentebox}{
  enhanced,breakable,pad at break=2mm,colback=milcGradeSoft,colframe=milcGradeDark,
  borderline west={3pt}{0pt}{milcGradeDark},
  arc=0mm,boxrule=0pt,left=5mm,right=4mm,top=2.5mm,bottom=2.5mm,
  before skip=4pt,after skip=8pt,
  fontupper=\itshape\small\color{milcNegro}\RaggedRight
}
% La flecha va en modo matematico a proposito: Helvetica Neue NO tiene el glifo
% U+2192, asi que \textrightarrow se compilaba a NADA («Missing character: There
% is no → in font Helvetica Neue») y el puente salia sin flecha. En modo
% matematico la flecha viene de la fuente matematica y si se dibuja.
\newcommand{\puente}[1]{%
  \begin{puentebox}%
    \textcolor{milcGradeDark}{$\rightarrow$}\hspace{2mm}#1%
  \end{puentebox}%
}

\newcommand{\linead}[1]{\noindent\color{milcLinea}\rule{#1}{0.5pt}\color{milcNegro}}
\newcommand{\respuesta}[1]{\par\vspace{2mm}\foreach \n in {1,...,#1}{\noindent\color{milcLinea}\rule{\linewidth}{0.5pt}\par\vspace{5mm}}\color{milcNegro}}
\newcommand{\checkbox}{\textcolor{milcGradeDark}{\fbox{\phantom{x}}}\hspace{5pt}}

% ─── Listas aireadas para las guías socioemocionales ────────────────────
% Componer las verificaciones como «\checkbox texto\\» deja el texto que salta
% de línea debajo del cuadrito y apelmaza el bloque. Estos entornos alinean la
% continuación y airean los ítems. Los usa Territorio Interior; cualquier
% familia puede hacerlo.
\newlist{checklista}{itemize}{1}
\setlist[checklista]{
  label=\textcolor{milcGradeDark}{\fbox{\phantom{x}}},
  leftmargin=9mm, labelsep=3.2mm, itemsep=2.4mm,
  topsep=2.5mm, parsep=0pt, partopsep=0pt, before=\RaggedRight
}
\newlist{pasoslista}{enumerate}{1}
\setlist[pasoslista]{
  label=\textcolor{milcGradeDark}{\sffamily\bfseries\arabic*.},
  leftmargin=8mm, labelsep=3mm, itemsep=2.8mm,
  topsep=1mm, parsep=0pt, partopsep=0pt, before=\RaggedRight
}

\newcommand{\phasecard}[4]{
\begin{tcolorbox}[enhanced,colback=#3,colframe=#2,arc=3mm,boxrule=.5pt,height=3.05cm,valign=center,halign=center,left=1.5mm,right=1.5mm,top=2mm,bottom=2mm]
{\sffamily\bfseries\fontsize{22}{24}\selectfont\textcolor{#2}{#1}}\par
{\sffamily\bfseries\small #4}
\end{tcolorbox}
}

% ── Piezas del rediseno 2026 (legibilidad 6.º-11.º) ───────────────────
% Banda de estacion: reemplaza el titlebox macizo. Titulo a la izquierda,
% dato practico (tiempo/modalidad) a la derecha, en una sola linea.
% Nunca huerfana: pide 9 lineas libres (o salta de pagina rellenando con
% \vfill) y prohibe el salto justo despues de la banda. Nueve y no seis:
% la caja partible que sigue necesita ~80pt (titulo adosado, relleno y dos
% lineas) para arrancar en la misma pagina; con menos, tcolorbox la manda
% entera a la siguiente y la banda se queda sola. El penalty 10000 va
% en `after`, ANTES del espacio posterior: un `after skip` (aunque sea 0pt)
% es un \vskip tras la caja, y ese glue es un punto de corte legal que un
% \nopagebreak escrito despues de \end{tcolorbox} ya no cierra. Cada banda
% deja marcador PDF y actualiza la cabecera derecha.
\newcounter{milcestacion}
\newcommand{\milcestacion}[2]{%
\Needspace*{9\baselineskip}%
\stepcounter{milcestacion}%
\pdfbookmark[1]{#2}{milcestacion\themilcestacion}%
\markright{#2}%
\begin{tcolorbox}[enhanced,colback=#1,colframe=#1,arc=2mm,boxrule=0pt,
  left=5mm,right=5mm,top=2.6mm,bottom=2.6mm,before skip=11pt,
  after={\par\nopagebreak\vskip7pt}]
{\sffamily\bfseries\fontsize{15}{18}\selectfont\milctintasobre{#1}#2}
\end{tcolorbox}}

% Tarjeta de paso numerada: sustituye las tablas «Pilar / Aplicacion», que
% eran formato de consulta docente, no de estudiante. El numero va en un
% circulo sobre el borde para que la secuencia se lea de un vistazo.
\newcommand{\milcpaso}[3]{%
\Needspace{4\baselineskip}%
\begin{tcolorbox}[enhanced,colback=white,colframe=milcLinea,arc=1.5mm,boxrule=.9pt,
  left=14mm,right=4mm,top=3mm,bottom=3mm,before skip=4pt,after skip=4pt,
  fontupper=\RaggedRight,
  overlay={\node[anchor=north west,circle,fill=#2,text=white,inner sep=0pt,
    minimum size=7.4mm,font=\sffamily\bfseries\small]
    at ([xshift=3.6mm,yshift=-3.2mm]frame.north west) {#1};}]
#3
\end{tcolorbox}}

% Casilla marcable de tamano util con lapiz (la anterior era \fbox{\phantom{x}},
% de alto variable segun el texto que la seguia).
\newcommand{\milcmarca}{\framebox[3.8mm]{\rule{0pt}{3.4mm}}}

% Fila marcable de lista suelta (checks de la estacion 1).
\newcommand{\milccheckrow}[1]{%
\par\vspace{1.8mm}\noindent
\makebox[6.4mm][l]{\raisebox{-.55mm}{\textcolor{milcNegro}{\milcmarca}}}%
\parbox[t]{\dimexpr\linewidth-6.4mm\relax}{\RaggedRight #1}\par}

% Celda marcable dentro de la rubrica (tabla de autoevaluacion). Misma
% sangria francesa, pero pensada para vivir dentro de una columna X.
\newcommand{\milcopcion}[1]{%
\makebox[5.2mm][l]{\raisebox{-.55mm}{\textcolor{milcNegro}{\milcmarca}}}%
\parbox[t]{\dimexpr\linewidth-5.2mm\relax}{\RaggedRight #1}}

% ── Recursos visuales de la guía ──────────────────────────────────────
% Declarados en el bloque `recursos:` del YAML y emitidos por el builder.
% \guiaFigura   → imagen rasterizada o PDF (foto, carta celeste, montaje).
% \guiaDiagrama → fuente TikZ/circuitikz (.tex) incrustada como vector:
%                 es la vía para circuitos y esquemáticos editables.
% [#1] = texto alternativo (alt) · #2 = ruta relativa al .tex · #3 = pie
% El alt llega al PDF etiquetado (/Alt) y lo lee el lector de pantalla.
\newcommand{\guiaFigura}[3][]{%
\begin{center}
\includegraphics[width=0.88\linewidth,height=0.38\textheight,keepaspectratio,alt={#1}]{#2}\\[1.5mm]
{\footnotesize\itshape\textcolor{milcNegro!75}{#3}}
\end{center}
}

\newcommand{\guiaDiagrama}[2]{%
\begin{center}
\input{#1}\\[1.5mm]
{\footnotesize\itshape\textcolor{milcNegro!75}{#2}}
\end{center}
}

\begin{document}
\thispagestyle{empty}

% ── Portada ───────────────────────────────────────────────────────────
% Antes: un rectangulo de color a pagina completa. Se veia bien en pantalla,
% pero en la fotocopiadora del colegio se comia el toner y salia gris sucio.
% Ahora la banda ocupa el tercio superior y el resto va en blanco.
\begin{tikzpicture}[remember picture,overlay]
  \path[fill=milcGradePrimary]
    (current page.north west) rectangle ([yshift=-10.6cm]current page.north east);

  \node[anchor=north west,text=milcGradeText] at ([xshift=2.0cm,yshift=-1.55cm]current page.north west)
    {\sffamily\bfseries\fontsize{12}{14}\selectfont GRADO};
  \node[anchor=north west,text=milcGradeText,opacity=.85] at ([xshift=2.0cm,yshift=-2.35cm]current page.north west)
    {\sffamily\bfseries\fontsize{8.5}{10}\selectfont SERIE GUÍAS · TECNOLOGÍA E INFORMÁTICA};
  \node[anchor=north east,text=milcGradeText,opacity=.85,align=right] at ([xshift=-2.0cm,yshift=-1.55cm]current page.north east)
    {\sffamily\bfseries\fontsize{9}{12}\selectfont I.E. Sor María Juliana\\Cartago, Valle del Cauca};

  \node[anchor=west,text=milcGradeText] (gradonum) at ([xshift=1.85cm,yshift=-7.1cm]current page.north west)
    {\sffamily\bfseries\fontsize{112}{112}\selectfont 8};
  % El circulito de grado (°) solo aplica a las guias de grado («11°»); en
  % Bebras y Semillero el numero grande es un momento/modulo, no un grado.
  % Se posiciona relativo al nodo `gradonum`, no en coordenadas absolutas.
  \draw[milcGradeText,line width=1.6pt]
    ([xshift=2.0mm,yshift=-4.5mm]gradonum.north east) circle (.15cm);

  \node[anchor=west,text=milcGradeText,text width=11.4cm,align=left] at ([xshift=1.5cm,yshift=0.5cm]gradonum.east)
    {
     {\sffamily\bfseries\fontsize{9}{11}\selectfont\MakeUppercase{Período 3 · Multimedia, narrativa y ciberseguridad}}\\[3mm]
     {\sffamily\bfseries\fontsize{21}{25}\selectfont\setlength{\baselineskip}{27pt}Sustentación pública\\[4pt]y portafolio digital\\[4pt]del grado 8}\\[3.5mm]
     {\sffamily\bfseries\fontsize{15}{18}\selectfont Guía 30-8-TIC}
    };
\end{tikzpicture}

\vspace*{9.4cm}
\vfill

{\sffamily\bfseries\fontsize{8}{10}\selectfont\textcolor{milcGradeDark}{LO QUE TE LLEVAS HOY}}

\begin{tcolorbox}[enhanced,colback=white,colframe=milcNegro,arc=2mm,boxrule=1.6pt,
  left=5mm,right=5mm,top=4mm,bottom=4mm,before skip=3pt,after skip=12pt,
  fontupper=\RaggedRight\fontsize{11}{15.5}\selectfont]
El cierre del año en tres piezas: un portafolio digital con seis trabajos curados de los tres periodos, cada uno con dos o tres líneas que digan qué demuestra; una carta abierta de cuatro párrafos, con algo concreto aprendido en cada periodo y un compromiso para noveno, firmada y fechada; y la sustentación pública de diez minutos, con demostración en vivo del proyecto integrador y cinco minutos de preguntas del grupo.
\end{tcolorbox}

\begin{tcolorbox}[enhanced,colback=milcGris,colframe=milcGris,arc=2mm,boxrule=0pt,
  left=5mm,right=5mm,top=3.5mm,bottom=3.5mm,after skip=12pt]
{\sffamily\bfseries\fontsize{8}{10}\selectfont\textcolor{milcNegro!55}{ESTA GUÍA ES DE}}\\[3mm]
\begin{tabularx}{\linewidth}{X p{3.2cm} p{3.2cm}}
\linead{\linewidth} & \linead{3.0cm} & \linead{3.0cm}\\[-1mm]
{\fontsize{8.5}{10}\selectfont\textcolor{milcNegro!55}{Nombre completo}} &
{\fontsize{8.5}{10}\selectfont\textcolor{milcNegro!55}{Grupo}} &
{\fontsize{8.5}{10}\selectfont\textcolor{milcNegro!55}{Fecha}}\\
\end{tabularx}
\end{tcolorbox}

\vfill

\noindent\textcolor{milcLinea}{\rule{\linewidth}{1pt}}\\[2mm]
{\sffamily\bfseries\fontsize{9.5}{12}\selectfont Autor: Dr. Álvaro Cárdenas Orozco}\\[1mm]
{\sffamily\fontsize{8.5}{11}\selectfont\textcolor{milcNegro!55}{Metodología MILC · apertura ancestral · pensamiento computacional · triángulo Dussel–estoicismo–Floridi}}

\newpage

\section*{Apertura: Sustentación pública y portafolio digital del grado 8}

\begin{softbox}{milcGradeDark}{milcGradeSoft}{Saber ancestral}
En los pueblos indígenas del Cauca, la autoridad no se toma: se recibe y se devuelve. El \textbf{cabildo} es elegido y reconocido por la propia comunidad, y su periodo dura un año. Así lo define el Decreto 2164 de 1995: una entidad pública especial, «elegidos y reconocidos por ésta». A quien recibe el cargo se le entrega un bastón de mando de madera de chonta con cintas de colores. Lo carga durante ese año y después pasa a otra persona. Desde 1991, el artículo 246 de la Constitución reconoce que esas autoridades pueden juzgar dentro de su territorio con sus propias normas. Lo que enseña esa figura es preciso: el mandato es temporal y termina rindiendo cuentas. La \textbf{cara de exclusión} está en el origen legal. La ley de 1890 que reconoció los cabildos lleva un título explícitamente racista, y la Corte Constitucional ya derribó tres de sus artículos. Hoy te toca a ti devolver lo que recibiste durante un año, delante de tu grupo.\par\smallskip{\footnotesize\textcolor{milcNegro!70}{\textbf{Fuente.} Presidencia de la República de Colombia. (1995, 7 de diciembre). \emph{Decreto 2164 de 1995}.}}
\end{softbox}

\begin{softbox}{milcTurquesa}{milcGris}{Saber contemporáneo}
Una sustentación de diez minutos no es un resumen del año: es una selección. Tiene cuatro bloques con tiempo asignado. \textbf{Proyecto integrador} (3 min): qué hace, la voz de la persona afectada y la propuesta. \textbf{Recorrido del portafolio} (3 min): cuatro o cinco piezas y qué demuestra cada una. \textbf{Carta abierta} (2 min): se lee completa. \textbf{Cierre} (2 min): qué te llevas a noveno. Después vienen cinco minutos de preguntas. El \textbf{portafolio digital} puede ser una carpeta compartida, un sitio sencillo o un blog. El formato da igual; lo que importa son tres cosas. Que cualquiera con el enlace pueda verlo. Que esté organizado por periodos. Que cada pieza tenga dos o tres líneas explicando qué demuestra. Un archivo suelto sin descripción no comunica nada a quien no estuvo en la clase.
\end{softbox}

\begin{softbox}{milcMostaza}{milcGris}{Pregunta-puente}
\textit{El cabildo devuelve el bastón al año y rinde cuentas de lo que hizo con él. Tú vas a devolver un año de trabajo en diez minutos. ¿Qué seis piezas cuentan de verdad lo que aprendiste?}
\end{softbox}

\begin{softbox}{milcVerde}{milcGris}{Saber y saber hacer}
Al terminar podrás: (1) \textbf{analizar} tu producción del año y curar seis piezas que demuestren aprendizaje, no facilidad; (2) \textbf{explicar} la estructura de la sustentación, del portafolio y de la carta, con su tiempo asignado; (3) \textbf{crear} el portafolio digital, la carta abierta firmada y la sustentación de diez minutos.
\end{softbox}



\small
\begin{tabularx}{\textwidth}{>{\bfseries}p{3.0cm}Y}
\rowcolor{milcGradePrimary!12}Autor & Dr. Álvaro Cárdenas Orozco\\
DBA & Producción multimedia ética y comportamiento responsable en entornos digitales (MEN, T\&I)\\
Referentes & Dussel (estética de la liberación) · Ley 1336 · Floridi (infoesfera)\\
Producto final & El cierre del año en tres piezas: un portafolio digital con seis trabajos curados de los tres periodos, cada uno con dos o tres líneas que digan qué demuestra; una carta abierta de cuatro párrafos, con algo concreto aprendido en cada periodo y un compromiso para noveno, firmada y fechada; y la sustentación pública de diez minutos, con demostración en vivo del proyecto integrador y cinco minutos de preguntas del grupo.\\
\end{tabularx}
\normalsize

\puente{El año empieza con datos, sigue con lógica y sensores y termina con multimedia y ética. En la sesión 9 juntaste todo en un proyecto sobre un problema real. Hoy lo sustentas y armas el portafolio con el que llegas a noveno.}

\milcestacion{milcGradeDark}{Ruta MILC: así vamos a aprender}
\begin{tabularx}{\textwidth}{>{\centering\arraybackslash}X>{\centering\arraybackslash}X>{\centering\arraybackslash}X>{\centering\arraybackslash}X}
\phasecard{1}{milcVerde}{milcGris}{Escucha\\[-1mm]\small Observo, escucho y pregunto.} &
\phasecard{2}{milcTurquesa}{milcGris}{Entender\\[-1mm]\small Sistematizo y ordeno.} &
\phasecard{3}{milcMagenta}{milcGris}{Hacer\\[-1mm]\small Construyo y aplico.} &
\phasecard{4}{milcVino}{milcGris}{Cierre\\[-1mm]\small Evalúo y reflexiono.}
\end{tabularx}


\puente{Antes de la teoría vas a revisar tu propia producción. Cuadernos, archivos, proyectos de los tres periodos. De ahí salen las seis piezas del portafolio, y elegir es la parte difícil.}

\milcestacion{milcVerde}{Estación 1 de 3 · Escucha}

\begin{softbox}{milcVerde}{milcGris}{Escena para escuchar}
\textbf{Actividad 1 · ANALIZA --- Curaduría del año} (15 min · individual). (1) Revisa tus cuadernos y tus archivos de los tres periodos. (2) Elige máximo seis piezas: dos por periodo, más el proyecto integrador. (3) Elige las que mejor muestren lo que aprendiste, no las que te salieron fáciles. (4) Escribe al lado de cada una, en una línea, qué demuestra. (5) Marca cuál pondrías primera si solo pudieras mostrar una.
\end{softbox}

{\sffamily\bfseries\fontsize{8}{10}\selectfont\textcolor{milcNegro!55}{MARCA CUANDO LO HAYAS HECHO}}
\milccheckrow{Revisé piezas de los tres periodos.}
\milccheckrow{Elegí máximo seis y escribí qué demuestra cada una.}
\milccheckrow{Marqué cuál pondría primera si solo pudiera mostrar una.}

\begin{infoband}{milcVerde}


\textbf{Lo que va al cuaderno (Actividad 1 · ANALIZA).} \textbf{Título:} «Actividad 1 --- Curaduría del año». \textbf{Formato:} lista de seis piezas con una línea cada una sobre qué demuestra, y la marca en la que iría primera. \textbf{Extensión:} un tercio de página. \textbf{Sabes que terminaste cuando} ninguna de las seis está por ser la más fácil.
\end{infoband}

\puente{Ya tienes la selección. Ahora vas a montar el plan minuto a minuto de los diez, el esqueleto del portafolio y la plantilla de la carta.}

\milcestacion{milcTurquesa}{Estación 2 de 3 · Entender}

\begin{softbox}{milcTurquesa}{milcGris}{Lo que vas a entender}
\textbf{Actividad 2 · EXPLICA --- El plan de los diez minutos} (30 min · parejas). (1) Con tu pareja, escriban el plan minuto a minuto de los cuatro bloques. (2) Armen el esqueleto del portafolio: los tres periodos y las seis piezas en su sitio. (3) Escriban la primera frase de cada párrafo de la carta, la que dice qué aprendieron en concreto. (4) Explíquenle su plan al otro en tres minutos y anoten qué no se entendió.
\end{softbox}

\milcpaso{1}{milcTurquesa}{\textbf{Los cuatro bloques, con su tiempo.} Proyecto integrador, tres minutos: qué hace, la voz de la persona y la propuesta. Recorrido del portafolio, tres minutos: cuatro o cinco piezas. Carta abierta, dos minutos: se lee completa. Cierre, dos minutos: qué te llevas a noveno. Después, cinco minutos de preguntas.}
\milcpaso{2}{milcTurquesa}{\textbf{El portafolio se juzga por las descripciones.} Un archivo suelto no dice nada a quien no estuvo en clase. Dos o tres líneas por pieza, con qué problema resolvía y qué aprendiste haciéndola. Eso es lo que convierte una carpeta en un portafolio.}
\milcpaso{3}{milcTurquesa}{\textbf{La carta, en cuatro párrafos.} Uno por periodo, cada uno con una cosa concreta, y el cuarto con un compromiso para noveno. Cierra con tu nombre completo y la fecha. Un párrafo bueno dice qué pasó: «el promedio del grupo era cuatro horas, pero un compañero tenía once». Uno malo dice «aprendí cosas sobre datos».}
\milcpaso{4}{milcTurquesa}{\textbf{Orgullo honesto, no falsa modestia.} Di con orgullo lo que hiciste bien y con honestidad lo que quedó a medias. «Mi micro:bit todavía falla cuando cambia la luz» y «armé un proyecto sobre el semáforo de mi cuadra» caben en la misma sustentación.}

\begin{drawbox}{milcTurquesa}{El cierre del grado 8, en una mirada}
\textbf{Sustentación (10 min):} proyecto 3 / portafolio 3 / carta 2 / cierre 2. \\[1mm] \textbf{Preguntas:} 5 minutos con honestidad. \\[1mm] \textbf{Portafolio:} 6 piezas con 2 o 3 líneas cada una. \\[1mm] \textbf{Carta:} 4 párrafos, nombre y fecha.
\end{drawbox}

\begin{softbox}{milcMostaza}{milcGris}{Errores comunes y su corrección}
(1) \textbf{Leer las diapositivas}: si se leen, sobra quien sustenta. (2) \textbf{Portafolio de archivos sueltos}: sin descripción, nadie sabe qué está viendo. (3) \textbf{Carta genérica}: «aprendí mucho» no es un párrafo, es un relleno. (4) \textbf{Esquivar la pregunta difícil}: se nota más esquivarla que no saber. (5) \textbf{Mostrar todo el año}: seleccionar seis piezas es parte de lo que se evalúa.
\end{softbox}

\begin{infoband}{milcTurquesa}


\textbf{Lo que va al cuaderno (Actividad 2 · EXPLICA).} \textbf{Título:} «Actividad 2 --- El plan de los diez minutos». \textbf{Formato:} el plan minuto a minuto, el esqueleto del portafolio por periodos y la primera frase de cada párrafo de la carta. \textbf{Extensión:} media página. \textbf{Sabes que terminaste cuando} tu pareja te dijo qué parte de tu plan no se entendió.
\end{infoband}

\puente{Tienes el plan. Toca ejecutarlo: armar el portafolio, escribir la carta, ensayar con cronómetro y sustentar delante del grupo.}

\milcestacion{milcMagenta}{Estación 3 de 3 · Hacer}

\begin{softbox}{milcMagenta}{milcGris}{Cómo vas a construir tu producto}
\textbf{Actividad 3 · CREA --- Portafolio, carta y sustentación} (25 min · individual). (1) Arma el portafolio con las seis piezas y sus descripciones. (2) Escribe la carta abierta de cuatro párrafos, fírmala y féchala. (3) Ensaya con cronómetro por lo menos dos veces. (4) Sustenta delante del grupo, con la demostración en vivo del proyecto integrador. (5) Responde las preguntas y anota la que te costó más. La sesión alcanza para armar y ensayar; el portafolio se termina en casa.
\end{softbox}

\milcpaso{1}{milcMagenta}{\textbf{Tu entrega tiene tres piezas.} El enlace del portafolio, abierto para cualquiera que lo tenga. La carta abierta firmada y fechada. La sustentación hecha dentro de los diez minutos.}
\milcpaso{2}{milcMagenta}{\textbf{Ensayar con cronómetro cambia el resultado.} Casi todo el mundo se pasa la primera vez. La segunda ronda es la que ajusta qué se cae. Si algo hay que recortar, se recorta del recorrido del portafolio, nunca de la carta.}
\milcpaso{3}{milcMagenta}{\textbf{La demostración en vivo se prepara para que falle.} Ten la pieza descargada y una captura de respaldo. Si algo no carga, sigues hablando. Que la conexión falle no es tu error; quedarte en blanco, sí.}
\milcpaso{4}{milcMagenta}{\textbf{Las preguntas son la parte que más se recuerda.} Contesta lo que sabes, y donde no sepas, dilo y di qué averiguarías. Cambiar de opinión delante de quien te corrige no debilita la sustentación: la termina de sostener.}

\begin{drawbox}{milcMagenta}{Antes de sustentar}
\checkbox Portafolio armado y con el enlace abierto para quien lo tenga. \\[1mm] \checkbox Las seis piezas con dos o tres líneas cada una. \\[1mm] \checkbox Carta abierta de cuatro párrafos, firmada y fechada. \\[1mm] \checkbox Ensayo cronometrado por lo menos dos veces. \\[1mm] \checkbox Proyecto integrador listo para demostrar en vivo. \\[1mm] \checkbox Respaldo descargado por si falla la conexión. \\[1mm] \checkbox Sustentación dentro de los diez minutos.
\end{drawbox}

\begin{softbox}{milcMagenta}{milcGris}{Plantilla de trabajo}
\textbf{Proyecto (3 min).} \textit{Qué hace, la voz, la propuesta.} \\[1mm] \textbf{Portafolio (3 min).} \textit{Cuatro o cinco piezas y qué demuestran.} \\[1mm] \textbf{Carta (2 min).} \textit{Cuatro párrafos leídos completos.} \\[1mm] \textbf{Cierre (2 min).} \textit{Qué me llevo a noveno.} \\[1mm] \textbf{Preguntas (5 min).} \textit{Lo que sé y lo que averiguaría.}
\end{softbox}

\begin{infoband}{milcMagenta}


\textbf{Lo que va al cuaderno (Actividad 3 · CREA).} \textbf{Título:} «Actividad 3 --- Portafolio, carta y sustentación». \textbf{Formato:} el enlace del portafolio, la carta completa, el plan minuto a minuto y las preguntas que te hicieron con lo que respondiste. \textbf{Extensión:} media página. \textbf{Sabes que terminaste cuando} anotaste la pregunta que más te costó y qué averiguarías para responderla mejor.
\end{infoband}

\puente{Lo que entregas hoy: el enlace del portafolio, la carta firmada y la sustentación hecha. La prueba de calidad: alguien que no estuvo en tus clases entiende, con solo abrir el portafolio, qué aprendiste.}

\milcestacion{milcMostaza}{Producto de la sesión}
\textbf{Portafolio digital de seis piezas + carta abierta firmada + sustentación de diez minutos}.
\begin{softbox}{milcMostaza}{milcGris}{Criterios de calidad}
(1) El portafolio tiene seis piezas de los tres periodos, cada una con su descripción. (2) El enlace lo puede abrir cualquiera que lo tenga. (3) La carta tiene cuatro párrafos concretos, nombre y fecha. (4) La sustentación cabe en diez minutos e incluye demostración en vivo. (5) Las preguntas se respondieron diciendo también lo que no se sabía.
\end{softbox}

\puente{Antes de cerrar, mira el año desde cinco lados. Reconocer con orgullo lo que salió bien y con honestidad lo que quedó a medias son la misma virtud, no dos.}

\milcestacion{milcVino}{Evaluación liberadora · cinco dimensiones}

\begin{softbox}{milcVino}{milcGris}{Sentido de la evaluación}
Evaluar no es solo revisar una nota. Es reconocer qué aprendí, qué cambió en mí, qué emociones manejé, cómo me relaciono con la ciudad y la comunidad, y qué le dejaré a quien viene detrás.
\end{softbox}

\textbf{Mi reflexión liberadora en cinco dimensiones.} Completa cada frase iniciada:

\small
\begin{tabularx}{\textwidth}{>{\bfseries\RaggedRight\arraybackslash}p{3.4cm}Y>{\RaggedRight\arraybackslash}p{4.6cm}}
\rowcolor{milcVino}\color{white}Dimensión & \color{white}\bfseries Mi respuesta (frame starter) & \color{white}\bfseries Algo que descubrí\\
1. Personal & Lo que más cambió en mí fue \linead{6.5cm} & \linead{4.2cm}\\
2. Emocional & La emoción más fuerte fue \linead{2.5cm} y la regulé \linead{3cm} & \linead{4.2cm}\\
3. Ciudadana & En democracia esto importa porque \linead{6cm} & \linead{4.2cm}\\
4. Local (Cartago) & En mi barrio sería útil para \linead{6.5cm} & \linead{4.2cm}\\
5. Intergeneracional & A mi abuela/abuelo le diría \linead{6.5cm} & \linead{4.2cm}\\
\end{tabularx}
\normalsize

\begin{infoband}{milcVino}
\textbf{Para la próxima sustentación.} Vuelve a esta tabla en seis meses. Verás que las respuestas cambian — no porque lo aprendido cambie, sino porque tú cambias. La evaluación liberadora es una conversación contigo a través del tiempo.
\end{infoband}


\puente{Para terminar, tres ideas sobre elegir y rendir cuentas, contadas con palabras de esta guía: Dussel sobre saber perder tiempo para elegir, Marco Aurelio sobre aceptar la corrección, Floridi sobre la abundancia que produce olvido.}

\milcestacion{milcGradeDark}{Tres ideas para cerrar · Dussel, un estoico y Floridi}

\begin{softbox}{milcVino}{milcGris}{Enrique Dussel · lente del nosotros}
\textit{Como los temas son infinitos y el tiempo corto, hay que saber perder tiempo eligiendo cuáles importan.}

{\footnotesize\itshape\color{milcNegro!70}--- idea tomada de Enrique Dussel, contada con palabras de esta guía. No es una cita textual.}

Los quince minutos que dedicas a escoger seis piezas valen más que una hora armando el portafolio con todo. Elegir es el trabajo; lo demás es montaje.

\textbf{¿Qué dejé fuera del portafolio, y estoy seguro de por qué?}
\end{softbox}
\linead{\linewidth}\\[1mm]
\linead{\linewidth}

\begin{softbox}{milcMostaza}{milcGris}{Estoicismo · Marco Aurelio · cuidado interior}
\textit{Eres igual de libre para cambiar de parecer que para seguir el consejo de quien te corrige.}

{\footnotesize\itshape\color{milcNegro!70}--- idea tomada de Marco Aurelio, contada con palabras de esta guía. No es una cita textual.}

En los cinco minutos de preguntas alguien va a señalar algo que no habías visto. Aceptarlo en el momento no te resta: muestra que entendiste tu propio trabajo lo bastante como para revisarlo.

\textbf{¿Qué me dijeron hoy que me hizo cambiar de opinión sobre algo que ya daba por hecho?}
\end{softbox}
\linead{\linewidth}\\[1mm]
\linead{\linewidth}

\begin{softbox}{milcTurquesa}{milcGris}{Luciano Floridi y la Onlife Initiative · lente de la infoesfera}
\textit{La abundancia de información también produce sobrecarga, distracción y olvido.}

{\footnotesize\itshape\color{milcNegro!70}--- idea tomada de Luciano Floridi y la Onlife Initiative, contada con palabras de esta guía. No es una cita textual.}

Un portafolio con treinta archivos sueltos no se lee: se cierra. Seis piezas con su descripción se recuerdan. Guardar todo es la forma más segura de que no quede nada.

\textbf{¿Mi portafolio ayuda a recordar lo que hice, o solo acumula lo que hice?}
\end{softbox}
\linead{\linewidth}\\[1mm]
\linead{\linewidth}

\begin{infoband}{milcNegro}
\textbf{Nota para el docente.} Las tres ideas están contadas con palabras de esta guía a partir del pensamiento de cada autor; no son citas textuales. De 9.º en adelante se trabajan con la obra y su referencia.
\end{infoband}

\milcestacion{milcVino}{Mi autoevaluación · marca una casilla por fila}

{\small
\setlength{\extrarowheight}{2.2pt}
\begin{tabularx}{\textwidth}{>{\RaggedRight\arraybackslash\bfseries}p{3.0cm}YYY}
\rowcolor{milcVino}\color{white}\bfseries Criterio & \color{white}\bfseries Lo logré & \color{white}\bfseries En proceso & \color{white}\bfseries Necesito apoyo\\[.6mm]
Comprensión del tema & \milcopcion{Explico las ideas con mis palabras.} & \milcopcion{Reconozco las ideas, falta precisión.} & \milcopcion{Necesito repasar lo básico.}\\[.8mm]
\rowcolor{milcGris}Pensamiento computacional & \milcopcion{Usé los pilares al armar mi producto.} & \milcopcion{Usé algunos pilares.} & \milcopcion{Necesito releer la estación 2.}\\[.8mm]
Aplicación y producto & \milcopcion{Mi producto cumple los criterios.} & \milcopcion{Está armado, con ajustes pendientes.} & \milcopcion{Aún está en construcción.}\\[.8mm]
\rowcolor{milcGris}Reflexión 5 dimensiones & \milcopcion{Respondí cada una con honestidad.} & \milcopcion{Respondí algunas con detalle.} & \milcopcion{Mis respuestas son muy generales.}\\[.8mm]
Triángulo de pensamiento & \milcopcion{Conecté las 3 voces con mi trabajo.} & \milcopcion{Conecté al menos una.} & \milcopcion{No las relaciono con mi proceso.}\\[.8mm]
\end{tabularx}}

\vspace{3mm}
\textbf{Hoy entendí que:}\\[1mm]
\linead{\linewidth}\\[1mm]
\linead{\linewidth}

\textbf{Mi compromiso para la próxima sesión es:}\\[1mm]
\linead{\linewidth}\\[1mm]
\linead{\linewidth}

\begin{softbox}{milcMostaza}{milcGris}{Cita de despedida · Marco Aurelio}
\textit{``El obstáculo en el camino se vuelve el camino. Lo que estorba al camino se vuelve camino.''}\\[1mm]
Cada error de hoy, bien observado, se convierte en pieza del camino que pisarás mañana.
\end{softbox}

\small
\begin{tabularx}{\textwidth}{>{\bfseries}p{3.6cm}Y}
\rowcolor{milcGradeDark!12}Nombre completo: & \linead{10cm}\\
Fecha: & \linead{4cm} \quad Curso/grupo: \linead{4cm}\\
Mi firma: & \linead{9cm}\\
\end{tabularx}
\normalsize

\begin{infoband}{milcGradeDark}
\textbf{Para el año que viene.} Guarda el enlace del portafolio donde puedas encontrarlo en enero. La carta firmada es tuya, no del colegio. En noveno vas a trabajar con bases de datos y con análisis de información, y el proyecto de este año es el punto de partida.
\end{infoband}

\end{document}
